\section{Limitations and Future Work}
\label{sec:limitations}

\subsection{Limitations}
\begin{itemize}
  \item \textbf{Reference translations.} The KDE4 references appear to be machine-generated (Chapter~\ref{sec:setup}), so
        chrF measures closeness to them rather than to human translation. Comparisons between pipeline versions are
        unaffected because all versions use the same references.
  \item \textbf{Calculated baselines.} The previous pipeline's times for 1\% and 5\% commits, and the start-of-semester
        stage in Chapter~\ref{sec:trajectory}, are calculated from measured rates; each is labelled.
  \item \textbf{One machine type and short strings.} All timings are for an m5.large and for user-interface strings with
        a median length of two words. Longer text would change absolute times, not the conclusions about what dominates them.
  \item \textbf{Experiment size.} The memory experiment used a 2,000-string application, in which near-twins are rarer
        than in the full corpus (1.5\% against 5.2\%).
  \item \textbf{Metric coverage.} This report uses chrF and placeholder survival; BLEU, ROUGE-L, METEOR and COMET are
        computed by Track~1.
\end{itemize}

\subsection{Future Work}
\begin{itemize}
  \item Repeat the memory experiment at full scale (10,054 strings) to confirm the savings on a production-size application.
  \item Compare machine types, including newer CPUs and GPUs, by cost per 1,000 translations.
  \item Evaluate beam~1 for the languages where its loss was smallest (under 1.3 chrF for Hindi, Marathi, Punjabi and
        Bengali), keeping beam~2 for Chinese and Japanese.
  \item Store the memory as per-language files in S3 now that it is loaded once per run, removing the always-on database.
\end{itemize}
